\documentclass[10pt,a4paper]{amsart}
\usepackage[margin=19mm]{geometry}
\usepackage{amsmath,amssymb,mathtools}
\usepackage[scr=rsfs,cal=euler]{mathalfa}
\usepackage{xfrac}
\usepackage[unicode,hidelinks,bookmarks=false]{hyperref}
\newcommand{\ie}{i.e.\ }
\newcommand{\cf}{cf.\ }
\newcommand{\resp}{respectively\ }
\newcommand{\Subsection}{\subsection}
\providecommand{\nobreakdash}{\nobreak\hbox{-}}
\setlength{\emergencystretch}{2em}
\multlinegap0pt
\usepackage{amssymb}
\usepackage{mathtools}
\usepackage[shortlabels]{enumitem}
\newtheorem{lemma}{Lemma}
\newcommand{\Zn}{\mathbb Z_n}
\title{The Type III realisation conjecture of Kirkland and \v{S}migoc}
\author{Brecht Verbeken, Vincent Ginis}
\date{Source: arXiv:2607.27219v1 (CC BY 4.0)}
\begin{document}
\maketitle

\noindent\emph{Source and licence.} The Type III realisation conjecture of Kirkland and \v{S}migoc. Version arXiv:2607.27219v1 (CC BY 4.0), distributed by arXiv under the Creative Commons Attribution 4.0 International licence, http://creativecommons.org/licenses/by/4.0/, as recorded in the arXiv licence metadata for this exact version and shown on the abstract page https://arxiv.org/abs/2607.27219v1. Full licence notice: \texttt{licenses/arxiv-2607.27219-CC-BY-4.0.txt}.\par\smallskip

\noindent\textit{Editorial scope.} Counted unit: the article's lemma lem:telescoping (source0.tex lines 809--814). The article's complete original proof of it is delivered with it (source0.tex lines 816--958, one \texttt{proof} environment of 143 lines). No mathematics is written, repaired or re-derived: the statement and the proof are the article's own lines, in the article's own notation, and no retained line is substituted or re-flowed.  Retained uncounted as the source-local context the proof needs: lines 677--701.  Nothing was omitted from any retained span, so the package carries no omission record.  No retained line contains a citation, so the wrapper carries no bibliography and no bibliography entry.  No retained line contains a figure environment, an image-inclusion command or drawing code, so no image is delivered and the package declares no asset dependency.  Editorial formatting only, in addition: an AMS \texttt{amsart} preamble that restates the article's own declarations and macros used by the retained lines, namely the environments \texttt{lemma} on separate counters, together with the standard packages those lines need.  The retained environments print in this document as Lemma 1 in order of appearance, whereas the article prints the same items with the numbering of its own full document.  The complete native source of the article as distributed in the arXiv e-print for this exact version is bundled unchanged as \texttt{sources/206338/source0.tex}.
\begin{lemma}\label{lem:turan-step}
Let $G$ be the graph on vertex set $S$ defined by
\[
        \{i,j\}\in E(G)\Longleftrightarrow Q_i\cap Q_j=\varnothing.
\]
Equivalently, since $n=qd+y\ge2q+1$,
\begin{equation}
        \{i,j\}\in E(G) \Longleftrightarrow d_n(i,j)\ge q. \label{eq:dist}
\end{equation}
The graph $G$ is a complete $d$-partite graph
\[
        S=S_1\sqcup\cdots\sqcup S_d
\]
with
\[
        \sum_{i\in S_r}w_i=\beta,\qquad r=1,\ldots,d.
\]
Consequently,
\[
        i,j\in S_r \Longrightarrow d_n(i,j)<q,
\]
\[
        i\in S_r,\ j\in S_s,\ r\ne s \Longrightarrow d_n(i,j)\ge q.
\]
\end{lemma}

\begin{lemma}[Circular-arc telescoping]\label{lem:telescoping}
For each part $S_r$ in Lemma \ref{lem:turan-step},
\[
        \sum_{i\in S_r}w_i=1-\prod_{i\in S_r}u_i.
\]
\end{lemma}

\begin{proof}
Fix $r\in\{1,\ldots,d\}$. By Lemma \ref{lem:turan-step},
\[
        \sum_{i\in S_s} w_i=\beta>0
\]
for every $s$, and hence every part $S_s$ is nonempty. For each $s\ne r$,
choose one index
\[
        j_s\in S_s.
\]
Since vertices belonging to different parts of the complete $d$-partite graph
$G$ are adjacent, Lemma \ref{lem:turan-step} gives
\[
        d_n(j_s,j_t)\ge q\qquad (s\ne t),
\]
and therefore the intervals $Q_{j_s}$, $s\ne r$, are pairwise disjoint.
Moreover, again by Lemma \ref{lem:turan-step}, if $i\in S_r$ and $s\ne r$,
then
\[
        d_n(i,j_s)\ge q,
\]
so $Q_i\cap Q_{j_s}=\varnothing$. Thus every interval $Q_i$, $i\in S_r$, is
contained in the complement of the union
\[
        \bigcup_{s\ne r} Q_{j_s}.
\]

Remove the $d-1$ disjoint intervals $Q_{j_s}$, $s\ne r$, from the cyclic
vertex set $\Zn$. The complement has
\[
        n-(d-1)q=qd+y-(d-1)q=q+y
\]
vertices. Since the intervals $Q_i$, $i\in S_r$, are pairwise intersecting by
Lemma \ref{lem:turan-step}, they must all lie in a single connected component
of this complement; denote that component by $H$. Indeed, two intervals
contained in different connected components of the complement would be
disjoint.

Let $L=|H|$. Since $H$ contains at least one interval $Q_i$ of length $q$, we
have $L\ge q$. On the other hand, $H$ is contained in a complement of total
size $q+y$, so
\[
        L\le q+y<2q.
\]
Write
\[
        L=q+\ell,\qquad 0\le \ell\le y<q.
\]
Cut the cycle open along $H$, and label the vertices of $H$ linearly as
\[
        1,2,\ldots,L.
\]
After this cut, every interval $Q_i\subseteq H$, $i\in S_r$, is an ordinary
linear interval of $q$ consecutive vertices. If its right endpoint is $t$,
then necessarily
\[
        q\le t\le L,
\]
because the interval has length $q$ and must remain inside $H$. Hence the
possible right endpoints lie in the interval
\[
        \{q,q+1,\ldots,L\},
\]
which has length
\[
        L-q+1=\ell+1\le q.
\]

List the right endpoints corresponding to indices in $S_r$ in increasing
linear order:
\[
        t_1<t_2<\cdots<t_m.
\]
Let $i_k\in S_r$ be the original index whose interval $Q_{i_k}$ has right
endpoint $t_k$ in this linear coordinate system.

Fix $k\in\{1,\ldots,m\}$. The cycle $Q_{i_k}$ consists of the backward edge
from $i_k$ to $i_k+1-q$, followed by the forward path
\[
        i_k+1-q\to i_k+2-q\to\cdots\to i_k.
\]
Its weight is therefore
\[
        w_{i_k}=(1-u_{i_k})
        \prod_{a=i_k-q+1}^{i_k-1}u_a,
\]
where the product is taken over the vertices on this forward path.

We now identify which factors in this product can be different from $1$. By
definition, $u_a<1$ exactly when $a\in S$. Thus only indices $a\in S$ matter.

First suppose $h<k$. Since all right endpoints $t_h$ lie in
$\{q,\ldots,L\}$, we have
\[
        t_k-t_h\le L-q=\ell<q.
\]
Therefore
\[
        t_k-q+1\le t_h\le t_k-1,
\]
so the original vertex $i_h$ appears among the indices in the product
defining $w_{i_k}$.

Next suppose $h>k$. Then $t_h>t_k$, so $i_h$ does not appear in the forward
path from $i_k+1-q$ to $i_k$.

Finally, let $a\in S_s$ for some $s\ne r$. Then, by the between-part condition
from Lemma \ref{lem:turan-step},
\[
        Q_a\cap Q_{i_k}=\varnothing.
\]
Since $a\in Q_a$, it follows that $a\notin Q_{i_k}$, and hence $a$ does not
appear in the product defining $w_{i_k}$.

Consequently, the only indices $a\in S$ appearing in the product
\[
        \prod_{a=i_k-q+1}^{i_k-1}u_a
\]
are precisely
\[
        i_1,\ldots,i_{k-1}.
\]
All other factors in the product have $a\notin S$, hence $u_a=1$. Therefore
\[
        w_{i_k}=(1-u_{i_k})\prod_{h<k}u_{i_h}.
\]
Summing over $k=1,\ldots,m$, we obtain the telescoping identity
\[
\begin{aligned}
        \sum_{i\in S_r}w_i
        &=(1-u_{i_1})
          +u_{i_1}(1-u_{i_2})
          +\cdots
          +u_{i_1}\cdots u_{i_{m-1}}(1-u_{i_m})\\
        &=1-u_{i_1}u_{i_2}\cdots u_{i_m}.
\end{aligned}
\]
Since $S_r=\{i_1,\ldots,i_m\}$, this is exactly
\[
        \sum_{i\in S_r}w_i=1-\prod_{i\in S_r}u_i.
\]
This proves the lemma.
\end{proof}

\end{document}
